\documentclass[11pt,a4paper]{article}

% ─── Packages ───────────────────────────────────────────────────────────────
\usepackage[top=1in,bottom=1in,left=0.9in,right=0.9in]{geometry}
\usepackage{amsmath,amssymb}
\usepackage{graphicx}
\usepackage{booktabs}
\usepackage{hyperref}
\usepackage{xcolor}
\usepackage{microtype}
\usepackage{authblk}
\usepackage{cite}
\usepackage{tabularx}
\usepackage{multirow}
\usepackage{url}
\usepackage{enumitem}
\usepackage{caption}

\hypersetup{
  colorlinks=true,
  linkcolor=blue!70!black,
  citecolor=blue!70!black,
  urlcolor=blue!70!black,
}

% ─── Title ──────────────────────────────────────────────────────────────────
\title{\textbf{Fine-Tuning Whisper for American English\\
Air Traffic Control Speech Recognition:\\
A Data-Efficient Pipeline}}

\author[1]{Jeffrey Su}
\author[2]{Omar Haq}
\affil[1]{Undergraduate Student, Rice University\\ \small\texttt{js318@rice.edu}}
\affil[2]{Undergraduate Student, University of Michigan\\ \small\texttt{omarhaq@umich.edu}}
\affil[ ]{\small\url{https://github.com/jeffreysupergoodatcoding/liveatc-pipeline}\ $\cdot$\
\url{https://huggingface.co/jeffreysuu/whisper-atc-finetuned}}
\date{February 2026}

% ────────────────────────────────────────────────────────────────────────────
\begin{document}
\maketitle

% ─── Abstract ───────────────────────────────────────────────────────────────
\begin{abstract}
Automatic speech recognition (ASR) for air traffic control (ATC) presents
a persistent domain adaptation challenge: VHF radio channel degradation,
specialized vocabulary, and a near-total absence of American English training
corpora in the published literature.
The state-of-the-art open ATC ASR model---WhisperATC by van Doorn et
al.~\cite{vandoorn2024whisperatc}, fine-tuned on European ATCO2 and ATCOSIM
corpora---achieves 3.88\% word error rate (WER) on ATCOSIM (speaker-split
evaluation) but degrades to 30.3\% on American ATC transmissions, exposing a
systematic accent and phraseology mismatch.
We present a data-efficient fine-tuning pipeline that adapts \textit{Whisper
Large v3}~\cite{radford2023whisper} to American English ATC using only 55
manually transcribed clips recorded from three major US airports (KIAH, KJFK,
KSFO) via LiveATC.net.
Domain-matched audio preprocessing---a 300--3400\,Hz Butterworth bandpass
filter with EBU\,R128 loudness normalization---combined with five-fold
stochastic data augmentation addresses the limited corpus size.
Full fine-tuning with conservative hyperparameters achieves \textbf{13.7\%
WER}, a \textbf{54.8\% relative reduction} from the European-trained baseline,
using 370$\times$ fewer training clips than the most comparable prior study.
A secondary contribution is the characterization of a structural incompatibility
between the HuggingFace PEFT LoRA implementation and Whisper's log-mel
spectrogram encoder that prevents parameter-efficient fine-tuning without
modification of library internals.
All code, the fine-tuned model, and training notebooks are publicly available.
\end{abstract}

\noindent\textbf{Keywords:} air traffic control, automatic speech recognition,
Whisper, fine-tuning, data augmentation, American English, domain adaptation,
LoRA

% ─── 1. Introduction ────────────────────────────────────────────────────────
\section{Introduction}

Air traffic control is among the most safety-critical communication domains in
civil aviation. Controllers manage the separation and flow of aircraft
exclusively through standardized voice radio exchanges governed by ICAO
Doc\,4444 and FAA Order\,7110.65, where a single transcription error or
phraseology misinterpretation can contribute to catastrophic outcomes---as
evidenced by the 1977 Tenerife runway collision, in which radio
miscommunication between controllers and crew contributed to the deadliest
accident in aviation history, and the 2023 Watsonville mid-air collision,
in which inadequate radio situational awareness was a contributing factor.

Despite this criticality, ATC-specific ASR research has lagged behind
general-domain speech recognition, and the available training corpora exhibit
a pronounced geographic bias. The ATCO2 corpus~\cite{kocour2021atco2} and the
ATCOSIM simulation dataset~\cite{zuluagagomez2022atcosim}---the two dominant
public resources---consist exclusively of European controller speech. Van Doorn
et al.~\cite{vandoorn2024whisperatc} recently achieved state-of-the-art
performance on these benchmarks by fine-tuning Whisper Large\,v3 on the
combined ATCO2 and ATCOSIM corpora, reporting 3.88\% WER on ATCOSIM
(speaker-split) and 13.5\% WER on ATCO2, and releasing the first
high-accuracy open models for ATC speech recognition.
However, when applied without modification to American English ATC
transmissions, their model (\texttt{jlvdoorn/atco2-asr-atcosim}) yields
30.3\% WER---a 26-percentage-point degradation attributable to compound
mismatches in accent, number verbalization conventions (e.g., ``one five zero''
vs.\ ``150''), and procedural phraseology between European and American ATC
practice.

American English ATC is therefore systematically underserved by the current
research landscape. The FAA projects a deficit of over 3,000 certified
professional controllers by 2032~\cite{faashortage2023}, and accurate
low-latency ASR is a prerequisite for any AI system designed to augment human
controllers, automate transmission logging, or surface real-time safety alerts.

This paper makes the following contributions:
\begin{enumerate}[noitemsep]
  \item A complete end-to-end pipeline for collecting, annotating, augmenting,
        and fine-tuning on American ATC audio sourced from LiveATC.net
        (Section~\ref{sec:method}).
  \item A 54.8\% relative WER reduction on American ATC transmissions from 55
        manually labeled clips, demonstrating that targeted fine-tuning with
        domain-matched preprocessing substantially closes the
        European--American accent and convention gap
        (Section~\ref{sec:experiments}).
  \item A documented structural incompatibility between the HuggingFace PEFT
        LoRA implementation and Whisper's audio encoder that constitutes a
        practical obstacle for the broader speech adaptation community
        (Section~\ref{sec:lora}).
\end{enumerate}

% ─── 2. Related Work ────────────────────────────────────────────────────────
\section{Related Work}
\label{sec:related}

\subsection{ATC Automatic Speech Recognition}

Early ATC ASR systems employed HMM-based hybrids augmented with task-specific
pronunciation grammars~\cite{cordero2012atcasr}. The ATCO2
project~\cite{kocour2021atco2} produced a large corpus of European ATC audio
and trained end-to-end models achieving sub-20\% WER on European data,
establishing the benchmark used by most subsequent work.
Pellegrini et al.~\cite{pellegrini2021atc} demonstrated that pre-trained
speech encoders (wav2vec\,2.0) fine-tuned on small ATC corpora achieve
competitive WER, directly motivating the transfer-learning approach pursued
here. Zuluaga-Gomez et al.~\cite{zuluagagomez2022atcosim} released ATCosim,
a simulation-based corpus that complements ATCO2 with synthetic air traffic
scenarios.

The most directly related work is van Doorn et
al.~\cite{vandoorn2024whisperatc}, who fine-tuned Whisper Large\,v3 on the
combined ATCO2 and ATCOSIM corpora, achieving 3.88\% WER on ATCOSIM
(speaker-split) and 13.5\% on ATCO2, and releasing the first high-accuracy
open models for this domain. Their work
establishes both the current state of the art and the baseline for our
evaluation: applied zero-shot to American ATC audio, the
\texttt{jlvdoorn/atco2-asr-atcosim} model yields 30.3\% WER. To our
knowledge, no published work has performed controlled fine-tuning specifically
targeting American English ATC.

\subsection{Whisper and Domain Adaptation}

OpenAI's Whisper~\cite{radford2023whisper} is a sequence-to-sequence
encoder-decoder model trained on 680,000 hours of weakly supervised web audio.
Its multitask, multilingual design makes it a strong candidate for domain
adaptation. Gandhe et al.~\cite{gandhe2023whisperadapt} showed that fine-tuning
Whisper on as few as 100 hours of in-domain speech can halve WER, while van
Doorn et al.~\cite{vandoorn2024whisperatc} established that even smaller
corpora yield strong results when the target domain is acoustically coherent
and domain-matched preprocessing is applied. Replicating the spectral
characteristics of the target recording environment is a consistent theme
across successful ASR domain adaptations~\cite{park2019specaugment}.

\subsection{Parameter-Efficient Fine-Tuning}

Low-Rank Adaptation (LoRA)~\cite{hu2021lora} reduces the number of trainable
parameters by inserting low-rank decomposition matrices into transformer
weight matrices, enabling fine-tuning of billion-parameter models on constrained
hardware. It has been applied to language
models~\cite{dettmers2023qlora} and certain speech
encoders~\cite{kim2023loraspeech}. Section~\ref{sec:lora} documents why LoRA
cannot be straightforwardly applied to Whisper's encoder using the current
HuggingFace PEFT implementation.

% ─── 3. Methodology ─────────────────────────────────────────────────────────
\section{Methodology}
\label{sec:method}

\subsection{Data Collection}

Audio was recorded from LiveATC.net, a public streaming service that
rebroadcasts live VHF radio transmissions from airports worldwide.
Three high-traffic American facilities were selected to maximize phonetic
and procedural diversity:
\begin{itemize}[noitemsep]
  \item \textbf{KIAH} --- Houston George Bush Intercontinental (Tower,
        Approach, Ground)
  \item \textbf{KJFK} --- New York John F.\ Kennedy International (Tower,
        Ground)
  \item \textbf{KSFO} --- San Francisco International (Tower, Ground,
        Departure)
\end{itemize}

A Node.js pipeline records continuous streams in 10-minute chunks and
applies silence-based segmentation (500\,ms silence threshold,
$-40$\,dB gate) to produce individual transmission segments, retaining
those within a 2--20 second duration window. Segments are uploaded to
Supabase cloud storage and surfaced through a custom Next.js web interface
that enables human annotators to play each segment, verify or correct a
draft transcription produced by Deepgram's commercial ASR engine, and
approve the result. Annotators followed a digit-first transcription
convention matching American ATC practice: numbers, callsigns, altitudes,
headings, and frequencies were written as digit strings (e.g., \texttt{425},
\texttt{135.7}, \texttt{240}) rather than verbalized form (e.g.,
\textit{four two five}, \textit{one three five point seven}). This convention
directly defines the target output format the fine-tuned model learns and
accounts for the numeric notation differences observed in
Section~\ref{sec:experiments}. After review, 55 high-quality clips were
accepted as ground-truth training data. The pipeline is summarized in
Figure~\ref{fig:pipeline}.

\begin{figure}[h]
\centering
\small
\begin{verbatim}
LiveATC.net  ->  Record (10-min chunks)
              ->  Silence-based segmentation (500 ms / -40 dB)
              ->  Upload to Supabase
              ->  Human annotation (Next.js web UI)
              ->  Export: audio + metadata.jsonl
              ->  Preprocessing & augmentation
              ->  Fine-tuning (HuggingFace Seq2SeqTrainer)
\end{verbatim}
\caption{End-to-end data pipeline from raw radio stream to fine-tuned model.}
\label{fig:pipeline}
\end{figure}

\subsection{Audio Preprocessing}
\label{sec:preprocess}

ATC transmissions are spectrally shaped by VHF radio hardware, concentrating
the energy of speech in the 300--3400\,Hz telephony band. To replicate this
characteristic during training and suppress broadband recording artifacts
not present in real operational audio, each clip is processed through a
6th-order Butterworth bandpass filter with cutoffs at 300\,Hz and 3400\,Hz.
Following filtering, loudness normalization to $-16$\,LUFS (EBU\,R128
standard) ensures uniform energy levels across clips recorded at varying
gain settings, preventing the model from learning spurious level-dependent
cues.

\subsection{Data Augmentation}
\label{sec:augmentation}

With only 55 base clips (approximately 15 minutes of labeled audio),
overfitting is the primary training risk. We synthesize four augmented
variants per clip, yielding an effective training set of 275 samples, via
the stochastic transforms listed in Table~\ref{tab:augmentation}.

\begin{table}[h]
\centering
\caption{Data augmentation transforms and parameter ranges.}
\label{tab:augmentation}
\begin{tabular}{ll}
\toprule
\textbf{Transform} & \textbf{Parameter Range} \\
\midrule
Time stretch         & Speed factor $\in [0.85,\,1.15]$ \\
Pitch shift          & $\pm 4$ semitones \\
Gain variation       & $\pm 6$\,dB \\
White noise injection & 1--8\% amplitude ratio \\
\bottomrule
\end{tabular}
\end{table}

Crucially, the 300--3400\,Hz bandpass filter is re-applied \emph{after} each
augmentation step. This ensures that all training variants---regardless of
time stretch, pitch shift, or noise level---exhibit the same radio spectral
envelope, preventing the model from associating augmentation artifacts with
the target domain.

\subsection{Model and Training Configuration}
\label{sec:training}

We fine-tune \texttt{openai/whisper-large-v3}~\cite{radford2023whisper}, a
1.55\,B-parameter encoder-decoder transformer, using the HuggingFace
\texttt{Seq2SeqTrainer} on a single NVIDIA Tesla T4 GPU (15\,GB VRAM) in
Google Colab. Training hyperparameters are reported in
Table~\ref{tab:hyperparams}.

\begin{table}[h]
\centering
\caption{Training hyperparameters.}
\label{tab:hyperparams}
\begin{tabular}{ll}
\toprule
\textbf{Hyperparameter} & \textbf{Value} \\
\midrule
Base model                  & \texttt{openai/whisper-large-v3} \\
Learning rate               & $5 \times 10^{-6}$ \\
Training epochs             & 5 \\
Per-device batch size       & 4 \\
Gradient accumulation steps & 4 (effective batch: 16) \\
Warmup steps                & 50 \\
Weight decay                & 0.01 \\
Numerical precision         & FP16 \\
Evaluation strategy         & Per epoch \\
Training duration           & $\approx$25 minutes \\
\bottomrule
\end{tabular}
\end{table}

The learning rate of $5 \times 10^{-6}$---approximately half of typical Whisper
fine-tuning recommendations ($1 \times 10^{-5}$)---was chosen to minimise
catastrophic forgetting of the model's general speech priors while permitting
specialization to American ATC conventions. Weight decay ($\lambda = 0.01$)
and per-epoch evaluation of validation loss serve as additional
regularizers against the limited data regime.

\subsection{LoRA Incompatibility with Whisper's Encoder}
\label{sec:lora}

Low-Rank Adaptation was initially pursued as a parameter-efficient alternative
to full fine-tuning that would reduce trainable parameters from 1.55\,B to
approximately 5--50\,M, lowering both memory requirements and overfitting
risk. We attempted LoRA via the HuggingFace PEFT library~\cite{peft2023} but
encountered a structural incompatibility: \texttt{get\_peft\_model} injects
adapters assuming the primary input key is \texttt{input\_ids}, which is the
standard token embedding path for language models. Whisper's encoder, however,
receives \texttt{input\_features}---a 128-channel log-mel spectrogram tensor
of shape $(B, 128, 3000)$. This mismatch produces a shape error at the first
forward pass that cannot be resolved without modifying PEFT library internals
or forking the encoder.

At the time of this work, no stable PEFT integration existed for
encoder-decoder speech models with Whisper's feature-extractor
front-end~\cite{kim2023loraspeech}. We therefore proceeded with full
fine-tuning, accepting the higher memory footprint and compensating with
the augmentation and conservative hyperparameters described above. This
incompatibility is documented here as it is likely to affect researchers
pursuing parameter-efficient adaptation of Whisper or architecturally
similar audio encoder-decoder models.

% ─── 4. Experiments ─────────────────────────────────────────────────────────
\section{Experiments and Results}
\label{sec:experiments}

\subsection{Evaluation Protocol}

WER is computed using the \texttt{jiwer} library~\cite{jiwer2023} with
standard text normalization (lowercasing, punctuation removal). The held-out
evaluation set comprises 15\% of the 55 annotated clips (approximately 8
segments) withheld from training and augmentation. The baseline system is the
\texttt{jlvdoorn/atco2-asr-atcosim} model from van Doorn et
al.~\cite{vandoorn2024whisperatc}, applied zero-shot to the same evaluation
clips.

\subsection{Main Results}

\begin{table}[h]
\centering
\caption{WER on held-out American English ATC clips.}
\label{tab:mainresults}
\begin{tabular}{lcc}
\toprule
\textbf{System} & \textbf{WER (\%)} & \textbf{Rel.\ Reduction} \\
\midrule
WhisperATC~\cite{vandoorn2024whisperatc} (zero-shot on American ATC) & 30.3 & --- \\
\textbf{Ours (Whisper Large v3, full fine-tune)}  & \textbf{13.7} & \textbf{54.8\%} \\
\bottomrule
\end{tabular}
\end{table}

The fine-tuned model achieves a 16.6 percentage-point absolute improvement
over the European-trained baseline (Table~\ref{tab:mainresults}), reducing
WER from 30.3\% to 13.7\% on American ATC audio. This result is notable
given the extreme data constraint: roughly 15 minutes of labeled audio,
compared to the thousands of hours typically used in large-scale ASR
fine-tuning.

\subsection{Qualitative Analysis}

Table~\ref{tab:qualitative} contrasts baseline and fine-tuned transcriptions
on three representative held-out segments, illustrating systematic improvements
in American-convention formatting.

\begin{table}[h]
\centering
\caption{Qualitative transcription comparison on held-out American ATC segments.}
\label{tab:qualitative}
\begin{tabularx}{\linewidth}{lX}
\toprule
& \textbf{Transcription} \\
\midrule
\textbf{Reference}  & United 425 cleared for takeoff runway 28 left wind 240 at 12 \\
\textbf{Baseline}   & United four two five clear for take off runway twenty-eight left wind two forty at twelve \\
\textbf{Fine-tuned} & United 425 cleared for takeoff runway 28 left wind 240 at 12 \\
\midrule
\textbf{Reference}  & Delta 832 descend and maintain flight level 240 speed 280 knots \\
\textbf{Baseline}   & Delta eight thirty two descend maintain flight level two forty speed two eighty knots \\
\textbf{Fine-tuned} & Delta 832 descend and maintain flight level 240 speed 280 knots \\
\midrule
\textbf{Reference}  & Southwest 1503 contact Houston departure 135.7 good day \\
\textbf{Baseline}   & Southwest Fifteen Zero Three contact Houston departure one three five point seven good day \\
\textbf{Fine-tuned} & Southwest 1503 contact Houston departure 135.7 good day \\
\bottomrule
\end{tabularx}
\end{table}

Three systematic improvements emerge across examples:
\begin{enumerate}[noitemsep]
  \item \textbf{Numeric notation}: American ATC renders callsigns, altitudes,
        headings, and frequencies as digit strings (e.g., ``425'', ``240'',
        ``135.7'') rather than the verbalized form produced by the
        European-trained model. This distinction directly affects downstream
        parsing of ATC instructions.
  \item \textbf{Phraseology completeness}: The conjunction \textit{and} in
        climb and descent clearances (\textit{descend and maintain}) is
        consistently recovered by the fine-tuned model; the baseline omits
        it, deviating from FAA\,7110.65 prescribed phraseology.
  \item \textbf{Command-word accuracy}: The fine-tuned model correctly
        transcribes \textit{cleared for takeoff} rather than the baseline's
        \textit{clear for take off}, preserving the operative word
        (\textit{cleared}) that distinguishes a clearance from a conditional.
\end{enumerate}

\subsection{Data Efficiency in Context}

\begin{table}[h]
\centering
\caption{Data efficiency comparison across recent ATC ASR fine-tuning studies.}
\label{tab:comparison}
\begin{tabular}{lrrc}
\toprule
\textbf{System} & \textbf{Training clips} & \textbf{WER (\%)} & \textbf{Accent region} \\
\midrule
van Doorn et al.~\cite{vandoorn2024whisperatc} & $\sim$10,000+ & 3.88 & European \\
Wee et al.~\cite{wee2023singaporeatc}          & $\sim$20,000  & $\sim$8 & Singaporean \\
\textbf{Ours}                                  & \textbf{55}   & \textbf{13.7} & \textbf{American} \\
\bottomrule
\end{tabular}
\end{table}

Our system uses approximately 370$\times$ fewer clips than Wee et
al.~\cite{wee2023singaporeatc} while targeting a distinct, previously
underserved accent group (Table~\ref{tab:comparison}). The remaining WER
gap relative to the European and Singaporean systems reflects two compounding
factors: the evaluation set is small ($\approx$8 clips), inflating variance,
and generalization from 55 training examples to the full phonetic and
procedural diversity of American ATC is necessarily limited. Scaling the
corpus to 1,000 or more clips would substantially reduce both sources of error.

% ─── 5. Discussion ──────────────────────────────────────────────────────────
\section{Discussion}
\label{sec:discussion}

\subsection{Limitations}

The evaluation set of approximately 8 clips introduces high variance in WER
estimates; results should be interpreted as indicative of the approach's
viability rather than as definitive performance bounds. Training data
covers only three airports; generalization to other American facilities
(e.g., KLAX, KORD, KATL) with distinct regional accents and traffic
patterns is untested. Full fine-tuning of a 1.55\,B-parameter model on 55
clips risks catastrophic forgetting of Whisper's multilingual and
general-domain capabilities; no out-of-domain evaluation was performed.
The documented LoRA incompatibility (Section~\ref{sec:lora}) means that
parameter-efficient adaptation---which would partially mitigate this
risk---requires either an upstream PEFT library update or a custom
adapter implementation.

\subsection{Future Work}

The most direct extension of this work is to scale the labeled American ATC
corpus to 1,000 or more clips spanning at least 10 facilities, bringing the
data scale closer to that of van Doorn et al.~\cite{vandoorn2024whisperatc}
and enabling statistically reliable WER estimation. The recording
infrastructure, annotation interface, and augmentation pipeline described in
Section~\ref{sec:method} are designed to support this scaling without
architectural changes. Structured transcriptions produced by the ASR layer
additionally serve as input for downstream natural language understanding,
extracting callsigns, instruction types, altitudes, and headings as
preconditions for AI-assisted workload reduction.

\subsubsection{Toward RLHF-Based ATC Reasoning}

A natural extension of the ASR foundation is a decision-support layer trained
via Reinforcement Learning from Human Feedback
(RLHF)~\cite{ouyang2022rlhf,christiano2017rlhf}. In this framing, a
policy model receives a structured world state (derived from ASR output plus
radar and surveillance feeds) and generates candidate ATC instructions.
Certified controllers rank candidate instructions on dimensions of safety,
efficiency, and phraseology compliance; a reward model is fitted on these
preference pairs and used to fine-tune the policy via Proximal Policy
Optimization (PPO)~\cite{schulman2017ppo}.

RLHF is particularly well-suited to ATC reasoning because the task is not
purely deterministic---experienced controllers can legitimately disagree on
optimal sequencing strategies under complex traffic loads---making human
preference a meaningful training signal that cannot be captured by a fixed
rule set alone. A key implementation consideration is that preference
elicitation should present scenarios in the same structured format produced
by the NLU layer, shielding controller raters from RL internals.

\subsubsection{Toward RLVR with Formal Safety Constraints}

For safety-critical deployment, verifiable reward signals are preferable to
learned reward models, which can be miscalibrated or subject to reward
hacking~\cite{skalse2022reward}. Reinforcement Learning with Verifiable
Rewards (RLVR)~\cite{lightman2023rlvr,guo2025deepseekr1} trains a policy
against rewards that are automatically computable from first principles,
removing the dependence on human raters at every training step. For ATC,
candidate verifiable reward signals include:

\begin{itemize}[noitemsep]
  \item \textbf{Separation compliance}: Binary or continuous reward based on
        whether projected trajectories maintain FAA minimum separation
        (3\,nmi lateral or 1,000\,ft vertical in Class\,B/C airspace).
  \item \textbf{Conflict probability}: Estimated from physics-based
        trajectory prediction over a configurable lookahead horizon.
  \item \textbf{Runway incursion risk}: Derived from a surface movement model
        incorporating aircraft position, taxi clearances, and hold-short
        instructions.
  \item \textbf{Phraseology correctness}: Verified against a formal grammar
        derived from ICAO Doc\,4444 and FAA Order\,7110.65.
  \item \textbf{Controller workload proxy}: Estimated from communication
        density, pending acknowledgements, and instruction completion latency.
\end{itemize}

\subsubsection{Simulation Environment}

Training and evaluating RLHF and RLVR policies requires a high-fidelity
simulator that can generate diverse, controllable ATC scenarios at scale.
BlueSky~\cite{hoekstra2016bluesky}, an open-source Python ATC simulator
developed at TU Delft, provides an appropriate foundation. A production
training environment would extend BlueSky with the following state variables
and sensor models:

\begin{itemize}[noitemsep]
  \item \textbf{Surveillance}: ADS-B transponder data (position, altitude,
        velocity, intent), primary and secondary radar returns, TCAS resolution
        advisory states.
  \item \textbf{Weather}: Wind fields (shear layers, gusts), convective cells,
        turbulence zones, precipitation, runway visual range, and ceiling.
  \item \textbf{Aircraft performance}: Type-specific speed envelopes,
        climb and descent performance tables, turn radii, and wake turbulence
        categories (Light, Medium, Heavy, Super under ICAO classification).
  \item \textbf{Communication}: Tower, ground, approach, and departure
        frequency management; ATIS broadcast state; radio-failure protocols.
  \item \textbf{Airspace}: Sector boundaries, published procedures (SIDs,
        STARs, approaches), temporary flight restrictions, and NOTAMs.
\end{itemize}

This simulation environment would allow RLVR training to proceed without
human annotation at each step, with the separation compliance and conflict
probability rewards computed deterministically from the simulator state.
The principal challenge is ensuring that scenarios sampled during training
adequately cover edge cases---emergency declarations, pilot deviations,
equipment failures---where formal specification of ``correct'' controller
behavior is non-trivial.

% ─── 6. Conclusion ──────────────────────────────────────────────────────────
\section{Conclusion}

We have demonstrated that Whisper Large\,v3 can be effectively adapted to
American English ATC communications using only 55 manually transcribed clips,
achieving 13.7\% WER---a 54.8\% relative improvement over the European-trained
WhisperATC baseline of van Doorn et al.~\cite{vandoorn2024whisperatc}. The
principal enabling factors are: (a) domain-matched audio preprocessing
(300--3400\,Hz bandpass filter re-applied after each augmentation step),
(b) five-fold stochastic augmentation that expands the effective training set
from 55 to 275 samples, and (c) conservative full fine-tuning hyperparameters
that minimize catastrophic forgetting.

A secondary contribution is the characterization of a structural
incompatibility between the HuggingFace PEFT LoRA implementation and
Whisper's log-mel spectrogram encoder. This finding is relevant to any
researcher pursuing adapter-based adaptation of encoder-decoder speech models
and points to an open engineering gap in the parameter-efficient fine-tuning
ecosystem.

The results demonstrate that targeted fine-tuning with domain-matched
preprocessing can substantially close accent-driven performance gaps even in
severe data-constrained settings, providing a data-efficient path toward
American English ATC speech recognition and a foundation for further research
into AI-assisted ATC operations.

% ─── Acknowledgements ───────────────────────────────────────────────────────
\section*{Acknowledgements}

The authors thank the operators of LiveATC.net for providing public access to
ATC audio streams, and Google for access to Colab GPU compute resources.

% ─── References ─────────────────────────────────────────────────────────────
\bibliographystyle{plain}
\bibliography{references}

\end{document}
